% Generated by scripts/assemble.py; corrections belong in source/corrections.json.
\examhead{53}{华东理工大学 2024 年数学分析试卷(回忆版)}{54}
% p055-b05; source PDF page 55
\par\Needspace{5\baselineskip}
\blockanchor{p055-b05}
\majorquestion{一}求极限.\par
% p055-b06; source PDF page 55
\blockanchor{p055-b06}
(1) $\lim_{n\to \infty}\sqrt[n]{2\sin^2{n} + \cos^2{n}}.$\par
% p055-b07; source PDF page 55
\blockanchor{p055-b07}
(2) $\lim_{x\to 0}\frac{\cos x - e^{-\frac{x^2}{2}}}{x^4}.$\par
% p055-b08; source PDF page 55
\par\Needspace{5\baselineskip}
\blockanchor{p055-b08}
\majorquestion{二}求重积分 $\iint_{D} e^{\max\{x^2, y^2\}} \mathrm{d}x \mathrm{~d}y$, 其中 $D: [0,1] \times [0,1]$.\par
% p055-b09; source PDF page 55
\par\Needspace{5\baselineskip}
\blockanchor{p055-b09}
\majorquestion{三}设 $f(x)$ 是 $[1, +\infty)$ 上单调递减的连续函数. 证明:\par
% p055-b10; source PDF page 55
\blockanchor{p055-b10}
(1) $\int_{1}^{+\infty} f(x) \, \mathrm{d}x$ 收敛时, 有 $\lim_{x \to +\infty} x f(x) = 0$.\par
% p055-b11; source PDF page 55
\blockanchor{p055-b11}
(2) 举例说明: 存在不单调递减的函数 $f(x)$ 使得 $\int_{1}^{+\infty} f(x) \, \mathrm{d}x$ 收敛, 但 $f(x)$ 在 $+\infty$ 的极限不存在.\par
% p055-b12; source PDF page 55
\par\Needspace{5\baselineskip}
\par\Needspace{6\baselineskip}
\blockanchor{p055-b12}
\majorquestion{四}求 $f(x, y, z) = \ln x + 2\ln y + 3\ln z$ 在 $x^2 + y^2 + z^2 = r^2 (x, y, z > 0, r > 0)$ 下的最大值，并证明对任意的正数 $a, b, c$，都有\par
\nopagebreak[4]
% p055-b13; source PDF page 55
\blockanchor{p055-b13}
\displayformula{a b ^ {2} c ^ {3} \leq 108 \left(\frac {a + b + c}{6}\right) ^ {6}.}
% p055-b14; source PDF page 55
\par\Needspace{5\baselineskip}
\par\Needspace{6\baselineskip}
\blockanchor{p055-b14}
\majorquestion{五}已知 $a_{n} > 0$, 证明:\par
\nopagebreak[4]
% p055-b15; source PDF page 55
\blockanchor{p055-b15}
\displayformula{\varliminf_ {n \to \infty} \frac {a _ {n + 1}}{a _ {n}} \leq \varliminf_ {n \to \infty} \sqrt [ n ]{a _ {n}} \leq \varlimsup_ {n \to \infty} \sqrt [ n ]{a _ {n}} \leq \varlimsup_ {n \to \infty} \frac {a _ {n + 1}}{a _ {n}}.}
% p055-b16; source PDF page 55
\par\Needspace{5\baselineskip}
\blockanchor{p055-b16}
\majorquestion{六}设函数列 $\{S_n(x)\}$ 在 $[0, a]$ 上连续, 且 $S_n(x) = \int_0^x S_{n-1}(t) \, \mathrm{d}t$, 证明: $\{S_n(x)\}$ 在 $[0, a]$ 上一致收敛.\par
% p055-b17; source PDF page 55
\par\Needspace{5\baselineskip}
\par\Needspace{6\baselineskip}
\blockanchor{p055-b17}
\majorquestion{七}求曲面积分\par
\nopagebreak[4]
% p055-b18; source PDF page 55
\blockanchor{p055-b18}
\displayformula{I = \iint_ {\Sigma} \frac {x \mathrm{d} y \mathrm{d} z + y \mathrm{d} z \mathrm{d} x + z \mathrm{d} x \mathrm{d} y}{\left(x ^ {2} + y ^ {2} + z ^ {2}\right) ^ {\frac {3}{2}}}.}
% p055-b19; source PDF page 55
\blockanchor{p055-b19}
其中  $\Sigma$  是立方体  $|x| \leq 1, |y| \leq 2, |z| \leq 3$  的表面, 取外侧.\par
% p055-b20; source PDF page 55
\par\Needspace{5\baselineskip}
\par\Needspace{6\baselineskip}
\blockanchor{p055-b20}
\majorquestion{八}证明函数\par
\nopagebreak[4]
% p055-b21; source PDF page 55
\blockanchor{p055-b21}
\displayformula{f (x) = \left\{ \begin{array}{l l} x ^ {\alpha} \cos \frac {1}{x}, & x \in \left(0, \frac {\pi}{2} \right]; \\ \left(x - \frac {\pi}{2}\right) ^ {2024} e ^ {- x}, & x \in \left(\frac {\pi}{2}, + \infty\right). \end{array} \right.}
% p055-b22; source PDF page 55
\blockanchor{p055-b22}
在 $(0, +\infty)$ 上一致连续, 其中 $\alpha > 0$.\par
% p055-b23; source PDF page 55
\par\Needspace{5\baselineskip}
\blockanchor{p055-b23}
\majorquestion{九}讨论 $\int_{1}^{+\infty}\cos (x^{p})\mathrm{d}x$ 的绝对收敛性和条件收敛性.\par
